\begin{afterword}
\summaryhead

A philosophical zombie is a being physically identical to you, atom for atom, but without consciousness. The zombie argument holds that if a whole zombie world is possible, consciousness is something beyond the physical. The post calls the resulting position epiphenomenalism: consciousness exists but has no effect on the physical world. It takes David Chalmers, author of \textsc{The Conscious Mind} (1996), as the argument's strongest advocate.

The post first argues that failing to see a contradiction in the zombie world does not show that the world is logically possible, much as failing to see a contradiction in a hard logic problem does not show it has a solution. It then argues that the zombie world is not even intuitive. When you attend to your own awareness, you think and say ``I am aware,'' which seems to be an effect of awareness. Yet zombie philosophers would write the same papers about consciousness, for the same physical reasons.

Chalmers accepts this, calling it a paradox, and answers that his beliefs are justified because he has the experiences, which his zombie twin lacks. The post replies that the physical part of Chalmers that writes his papers would then be right only by a separate coincidence, and that a self-inspecting AI would treat such a part as a fault. It compares three theories: interactionist substance dualism, a reductionism that has not yet explained consciousness, and epiphenomenal property dualism. It argues that the third posits the most unexplained ``miracles,'' and concludes that the sane response to the zombie argument is to say it cannot be right and look for the flaw.

\responsehead

The post's main objection is a serious and recognized one. That our talk about consciousness seems to be caused by consciousness had been pressed against epiphenomenalism before: by G.~F. Stout in 1931, and by Avshalom Elitzur in 1989, whom Chalmers credits in the book the post quotes. The post allows that the argument ``may well have been said before.'' Chalmers calls the problem ``the greatest tension that a nonreductive theory is faced with.'' In 2018 he quoted this post, cited it among ``coincidence arguments,'' replied to them, and conceded that ``it is hard to avoid a sense of coincidence entirely.'' The post's quotations from Chalmers are accurate, and its point that seeing no contradiction is not the same as logical possibility is one Chalmers makes himself.

Its account of its target is less accurate. The post says Chalmers ``specifies'' that consciousness is epiphenomenal. The book it quotes says ``I do not describe my view as epiphenomenalism,'' and leaves the causal question open, though it grants ``at least a weak form.'' By 2003 Chalmers gave credence to interactionism and panprotopsychism as well. The zombie argument by itself concludes only that physicalism is false. None of this defeats the main objection, since Chalmers says the paradox arises ``almost as strongly'' on the panprotopsychist view, and that interactionism is ``of little help'' against it.

The post's repeated question, why not let consciousness be causal, is answered in the book, and the post reports the answer only in part, running two of Chalmers's objections together. Chalmers cites evidence that physics is causally closed. His ``deepest reason'' is that the phenomenal can be subtracted from any causal story, including one with immaterial ``psychons,'' so interaction gains nothing. That also bears on the post's claim that substance dualism escapes the problem.

The closing comparison of theories rests on a premise Chalmers denies: that a reductionist account covers the same data. For Chalmers, experience is known directly, and he grants that without that evidence the simpler, zombie-like view would win. Whether experience is so known is the point in dispute. The post and Chalmers's replies (acquaintance, and phenomenal beliefs partly constituted by experience) leave it open; the post's reflective-coherence argument concerns the part of Chalmers an outsider can inspect, and Chalmers grants that such inspection cannot show whether he is conscious.

In short: a forceful statement of a recognized objection to epiphenomenalism, which Chalmers later called hard to answer fully, aimed at a Chalmers more committed to epiphenomenalism than the book it quotes, and without the book's main reason for not making consciousness causal.
\end{afterword}
